\documentclass[11pt, aspectratio=169]{beamer}

% ============================================================
% TEMA Y COLORES
% ============================================================
\usetheme{Madrid}
\usecolortheme{default}

\usepackage[utf8]{inputenc}
\usepackage[T1]{fontenc}
\usepackage[spanish,es-tabla]{babel}
\usepackage{booktabs}
\usepackage{graphicx}
\usepackage{listings}
\usepackage{xcolor}
\usepackage{amsmath,amssymb}
\usepackage{tabularx}
\usepackage{colortbl}
\usepackage{array}
\usepackage{tikz}
\usetikzlibrary{shapes.geometric, arrows.meta, positioning, fit, backgrounds}
\usepackage{tcolorbox}
\tcbuselibrary{skins}
\usepackage{pifont}
\usepackage{multirow}
\usepackage{lmodern}
\usepackage{fontawesome5}
\usepackage[nopatch=footnote]{microtype}

% ============================================================
% PALETA DE COLORES SMARTROOT
% ============================================================
\definecolor{SmartBlue}{RGB}{26,82,118}
\definecolor{SmartGreen}{RGB}{39,174,96}
\definecolor{SmartOrange}{RGB}{230,126,34}
\definecolor{SmartRed}{RGB}{192,57,43}
\definecolor{SmartGray}{RGB}{52,73,94}
\definecolor{SmartLightBlue}{RGB}{41,128,185}
\definecolor{SmartYellow}{RGB}{243,156,18}
\definecolor{backcolour}{rgb}{0.97,0.97,0.95}
\definecolor{codegreen}{rgb}{0,0.55,0}
\definecolor{codegray}{rgb}{0.5,0.5,0.5}
\definecolor{codepurple}{rgb}{0.58,0,0.82}

% ============================================================
% PERSONALIZACIÓN DEL TEMA
% ============================================================
\setbeamercolor{palette primary}{bg=SmartBlue,fg=white}
\setbeamercolor{palette secondary}{bg=SmartBlue!80,fg=white}
\setbeamercolor{palette tertiary}{bg=SmartBlue!60,fg=white}
\setbeamercolor{palette quaternary}{bg=SmartBlue!40,fg=white}
\setbeamercolor{title}{fg=SmartBlue}
\setbeamercolor{frametitle}{bg=SmartBlue,fg=white}
\setbeamercolor{structure}{fg=SmartBlue}
\setbeamercolor{block title}{bg=SmartBlue,fg=white}
\setbeamercolor{block body}{bg=SmartBlue!8,fg=SmartGray}
\setbeamercolor{block title alerted}{bg=SmartRed,fg=white}
\setbeamercolor{block body alerted}{bg=SmartRed!8,fg=SmartGray}
\setbeamercolor{block title example}{bg=SmartGreen,fg=white}
\setbeamercolor{block body example}{bg=SmartGreen!8,fg=SmartGray}
\setbeamerfont{frametitle}{size=\normalsize,series=\bfseries}
\setbeamerfont{title}{size=\LARGE,series=\bfseries}

% ============================================================
% ESTILO DE CÓDIGO
% ============================================================
\lstdefinestyle{smartstyle}{
    backgroundcolor=\color{backcolour},
    commentstyle=\color{codegreen}\itshape,
    keywordstyle=\color{SmartLightBlue}\bfseries,
    numberstyle=\tiny\color{codegray},
    stringstyle=\color{codepurple},
    basicstyle=\ttfamily\scriptsize,
    breakatwhitespace=false,
    breaklines=true,
    captionpos=b,
    keepspaces=true,
    numbers=left,
    numbersep=5pt,
    showspaces=false,
    showstringspaces=false,
    showtabs=false,
    tabsize=3,
    literate={á}{{\'a}}1 {é}{{\'e}}1 {í}{{\'i}}1 {ó}{{\'o}}1
             {ú}{{\'u}}1 {ñ}{{\~n}}1 {Ñ}{{\~N}}1,
    frame=single,
    rulecolor=\color{SmartGray!30},
    xleftmargin=12pt,
    xrightmargin=4pt
}
\lstset{style=smartstyle, language=Python}

% ============================================================
% COMANDOS AUXILIARES
% ============================================================
\newcommand{\semverde}[1]{\textcolor{SmartGreen}{\textbf{#1}}}
\newcommand{\semamarillo}[1]{\textcolor{SmartYellow}{\textbf{#1}}}
\newcommand{\semrojo}[1]{\textcolor{SmartRed}{\textbf{#1}}}
\newcommand{\tick}{\textcolor{SmartGreen}{\ding{52}}}
\newcommand{\kpirow}[5]{#1 & #2 & #3 & #4 & \cellcolor{SmartGreen!25}\semverde{#5} \\}

% ============================================================
% METADATOS
% ============================================================
\title[SmartRoot --- Decisión de Riego con Python]{%
  \textbf{SmartRoot}\\[4pt]
  {\large Sistema Inteligente de Decisión de Riego\\
  y Validación de Sensores Hídricos}\\[4pt]
  {\normalsize mediante Programación Científica en Python}}

\subtitle{Trabajo Final --- Fundamentos de Programación}

\author[Estudiante | Becerra Ph.D.]{%
  [Nombre del Estudiante]\\[2pt]
  {\small Docente: Miguel A. Becerra Ph.D.}}

\institute[Maestría ACI 2026-1]{%
  Maestría en Automatización y Control Industrial\\
  Fundamentos de Programación Científica · Posgrado 2026-1}

\date{5 de junio de 2026}

% ============================================================
% LOGO EN EL PIE
% ============================================================
% \logo{\includegraphics[height=0.6cm]{logo_universidad.png}}

% ============================================================
% DOCUMENTO
% ============================================================
\begin{document}

% ------------------------------------------------------------------
% DIAPOSITIVA 1: PORTADA
% ------------------------------------------------------------------
\begin{frame}[plain]
\begin{tikzpicture}[remember picture, overlay]
    % Fondo superior
    \fill[SmartBlue] (current page.north west) rectangle
        ([yshift=-0.38\paperheight]current page.north east);
    % Franja inferior
    \fill[SmartGreen!80!SmartBlue] (current page.south west) rectangle
        ([yshift=0.06\paperheight]current page.south east);
\end{tikzpicture}

\vspace*{0.3cm}
\begin{columns}[c]
\column{0.65\textwidth}
{\color{white}\LARGE\bfseries SmartRoot}\\[3pt]
{\color{white!85}\large Sistema Inteligente de Decisión de Riego}\\[2pt]
{\color{white!85}\large y Validación de Sensores Hídricos}\\[6pt]
{\color{white!70}\normalsize mediante Programación Científica en Python}\\[12pt]
\hrule width 0.9\textwidth height 0.3pt \relax
\vspace{6pt}
{\color{white!80}\small Trabajo Final · Fundamentos de Programación}

\column{0.35\textwidth}
\vspace{0.5cm}
\begin{tcolorbox}[colback=white!15!SmartBlue,colframe=white!40,
    arc=6pt, left=4pt, right=4pt, top=4pt, bottom=4pt]
    \centering
    {\color{white}\small\textbf{Dataset}}\\[2pt]
    {\color{white!85}\scriptsize Jackisch et al. (2020)}\\[1pt]
    {\color{white!85}\scriptsize JKI · Braunschweig, DE}\\[4pt]
    {\color{white}\small\textbf{2\,494 registros}}\\[1pt]
    {\color{white!85}\scriptsize Mayo--Julio 2016}\\[1pt]
    {\color{white!85}\scriptsize Res. 30 min}
\end{tcolorbox}
\end{columns}

\vspace{0.5cm}
\begin{columns}[c]
\column{0.5\textwidth}
{\small\textbf{Elaborado por:} [Nombre del Estudiante]}\\
{\small\textbf{Docente:} Miguel A. Becerra Ph.D.}

\column{0.5\textwidth}
{\small\textbf{Programa:} Maestría ACI}\\
{\small\textbf{Entrega:} 5 de junio de 2026}
\end{columns}
\end{frame}

% ------------------------------------------------------------------
% DIAPOSITIVA 2: TABLA DE CONTENIDO
% ------------------------------------------------------------------
\begin{frame}{Agenda de la Presentación}
\begin{columns}[t]
\column{0.5\textwidth}
\tableofcontents[sections={1-5}]
\column{0.5\textwidth}
\tableofcontents[sections={6-10}]
\end{columns}
\vspace{0.4cm}
\begin{tcolorbox}[colback=SmartBlue!8,colframe=SmartBlue!40,
    arc=4pt,left=6pt,right=6pt,top=3pt,bottom=3pt]
\centering
{\small Presentación de \textbf{15--20 minutos} · Sustentación técnica del proyecto}
\end{tcolorbox}
\end{frame}

\setcounter{tocdepth}{1}

% ==================================================================
\section{Problema y Contexto}
% ==================================================================

% ------------------------------------------------------------------
% DIAPOSITIVA 3: EL PROBLEMA
% ------------------------------------------------------------------
\begin{frame}{1a. Problema y Contexto}
\begin{alertblock}{Problema Central}
\begin{itemize}
    \item[\textbullet] Parcelas experimentales con \textbf{múltiples tecnologías} de sensores generan lecturas inconsistentes.
    \item[\textbullet] Sin validación cruzada automática, los errores de sensor se confunden con estrés hídrico real.
    \item[\textbullet] Decisiones de riego basadas en umbrales \textbf{determinísticos fijos} ignoran la incertidumbre del suelo.
\end{itemize}
\end{alertblock}

\vfill

\begin{block}{Pregunta de Investigación}
\centering
¿Cómo determinar automáticamente \textbf{cuándo, dónde y por qué} regar,\\
cuantificando el riesgo con base en datos reales de sensores?
\end{block}
\end{frame}

\begin{frame}{1b. Solución Propuesta (SmartRoot)}
\begin{block}{SmartRoot como solución integral}
\begin{itemize}
    \item[\tick] \textbf{Validación cruzada} entre tensiómetro (referencia) y Gypsum (secundario).
    \item[\tick] \textbf{5 KPI} con clasificación semáforo automática.
    \item[\tick] \textbf{Montecarlo} para cuantificar probabilidad de riego urgente.
    \item[\tick] \textbf{POO} para modelar sensores reales del dataset.
\end{itemize}
\end{block}

\vfill

\begin{tcolorbox}[colback=SmartGreen!10,colframe=SmartGreen!50,
    arc=4pt,left=6pt,right=6pt,top=4pt,bottom=4pt]
\centering
\textbf{Resultado global:} Los 5 KPI en estado \semverde{VERDE}\\
Manejo hídrico \textbf{excelente} en mayo--julio 2016
\end{tcolorbox}
\end{frame}

% ==================================================================
\section{Datos y Dataset}
% ==================================================================

% ------------------------------------------------------------------
% DIAPOSITIVA 4: DATASET
% ------------------------------------------------------------------
\begin{frame}{2a. Dataset Experimental}
\begin{block}{Descripción del dataset (Jackisch et al., 2020)}
\begin{itemize}
    \item \textbf{Fuente:} PANGAEA (doi:10.1594/PANGAEA.892319)
    \item \textbf{Parcela:} Julius Kühn-Institut (JKI), Braunschweig, Alemania
    \item \textbf{Período:} 13 mayo -- 4 julio 2016
    \item \textbf{Resolución:} 30 min por registro
    \item \textbf{Total validado:} 2\,494 registros
\end{itemize}
\end{block}

\vfill

\begin{tcolorbox}[colback=SmartOrange!12,colframe=SmartOrange!50,
    arc=4pt,left=6pt,right=6pt,top=4pt,bottom=4pt]
\textbf{Nota importante:} Los sensores Gypsum presentan una brecha de datos en los registros 1000--3600. Esto fue detectado y tratado automáticamente por el algoritmo de validación de SmartRoot.
\end{tcolorbox}
\end{frame}

\begin{frame}{2b. Sensores y Estadísticas}
\begin{columns}[c]
\column{0.45\textwidth}
\begin{exampleblock}{Red de Sensores}
\begin{itemize}
    \item \textbf{Tensióm.} T42--T54
    \item \textbf{Gypsum} Gyp1--4
    \item \textbf{FDR} ECTM1--4
    \item \textbf{Temp.} MPS6 1--4
\end{itemize}
\end{exampleblock}

\column{0.55\textwidth}
\begin{block}{Estadísticas (Media $\pm$ Std)}
\begin{itemize}
    \item $\Psi_{ref}$: $195.6 \pm 88.1$ hPa
    \item $\Psi_{gyp}$: $187.5 \pm 76.2$ hPa
    \item $\theta_{vol}$: $19.9 \pm 0.8$ \%
    \item $T_{suelo}$: $15.2 \pm 7.0$ °C
\end{itemize}
\end{block}
\end{columns}
\end{frame}

% ==================================================================
\section{Lógica del Programa}
% ==================================================================

% ------------------------------------------------------------------
% DIAPOSITIVA 5: LÓGICA - CARGA Y VALIDACIÓN
% ------------------------------------------------------------------
\begin{frame}[fragile]{3a. Lógica del Programa: Carga e Inspección}
\begin{block}{Carga e inspección inicial (Ciclo FOR)}
\begin{lstlisting}[numbers=none,frame=none,basicstyle=\ttfamily\small]
import pandas as pd
import numpy as np

# Carga del dataset corregido
df = pd.read_csv('data/processed/dataset_corregido.csv',
                 parse_dates=['Timestamp'],
                 index_col='Timestamp')

# Cobertura de sensores: ciclo for
cols_ref = ['T42','T43','T44','T51','T52','T54']
for col in cols_ref:
    cob = df[col].notna().mean() * 100
    print(f"{col}: {cob:.1f}%")
\end{lstlisting}
\end{block}
\end{frame}

\begin{frame}[fragile]{3b. Lógica del Programa: Validación e Imputación}
\begin{block}{Validación con condicionales múltiples y ciclo WHILE}
\begin{lstlisting}[numbers=none,frame=none,basicstyle=\ttfamily\small]
def validar_registro(psi, theta):
    errores = []
    # if / elif / else
    if psi is None or np.isnan(psi): errores.append("PSI_NULO")
    elif psi < 0: errores.append("PSI_NEGATIVO")
    elif psi > 1500: errores.append("FUERA_RANGO")
    
    return "VALIDO" if not errores else "INVALIDO"

# Ciclo while: imputacion iterativa temporal
it = 0
while df['psi_real'].isna().sum() > 0 and it < 5:
    df['psi_real'] = df['psi_real'].interpolate(method='time', limit=3)
    it += 1
\end{lstlisting}
\end{block}

\vfill
\begin{tcolorbox}[colback=SmartOrange!10,colframe=SmartOrange!50,
    arc=4pt,left=6pt,right=6pt,top=4pt,bottom=4pt]
\textbf{Resultado:} 2\,494 registros validados · Tasa de calidad $>98\%$ · Imputación en $\leq 3$ iteraciones
\end{tcolorbox}
\end{frame}

% ==================================================================
\section{Comprensiones Python}
% ==================================================================

% ------------------------------------------------------------------
% DIAPOSITIVA 6: LAS 3 COMPRENSIONES
% ------------------------------------------------------------------
\begin{frame}[fragile]{4a. Comprensiones Python (List y Dict)}
\begin{exampleblock}{\textbf{1.} List Comprehension --- Filtrado dinámico}
\begin{lstlisting}[numbers=none,frame=none,basicstyle=\ttfamily\small]
# Calcula media solo de sensores con datos > 95%
promedios_ref = [
    {'sensor': col, 'media': df[col].mean(), 
     'cobertura': df[col].notna().mean() * 100}
    for col in cols_psi_ref
    if df[col].notna().mean() > 0.95     # FILTRO INTELIGENTE
]  
# Resultado: Lista de 6 diccionarios validados.
\end{lstlisting}
\end{exampleblock}

\vfill
\begin{block}{\textbf{2.} Dictionary Comprehension --- Evaluación de Semáforo $O(1)$}
\begin{lstlisting}[numbers=none,frame=none,basicstyle=\ttfamily\small]
mapes = {'Gypsum1': 12.8, 'Gypsum2': 14.1, 'Gypsum3': 13.5}
# Mapeo directo: sensor -> estado
semaforo = {
    sid: ('VERDE' if m <= 20 else 'AMARILLO' if m <= 35 else 'ROJO')
    for sid, m in mapes.items()
}
# Resultado: {'Gypsum1': 'VERDE', 'Gypsum2': 'VERDE', ...}
\end{lstlisting}
\end{block}
\end{frame}

\begin{frame}[fragile]{4b. Comprensiones Python (Set)}
\begin{alertblock}{\textbf{3.} Set Comprehension --- Identificación de Estados Únicos}
\begin{lstlisting}[numbers=none,frame=none,basicstyle=\ttfamily\small]
# Extrae el conjunto matematico de todos los estados que realmente 
# ocurrieron en la serie temporal sin duplicados.
estados = {
    ('Humedo' if p<=100 else 'Optimo' if p<=300 else 
     'Seco' if p<=400 else 'Critico')
    for p in df['psi_real'].dropna()
}
# Resultado real en memoria: {'Humedo', 'Optimo', 'Seco'}
# Nota: La extraccion automatica demostro que "Critico" 
# jamas se registro.
\end{lstlisting}
\end{alertblock}
\end{frame}

% ==================================================================
\section{Indicadores KPI}
% ==================================================================

% ------------------------------------------------------------------
% DIAPOSITIVA 7: KPI COMPLETOS
% ------------------------------------------------------------------
\begin{frame}{3a. Indicadores de Desempeño (KPI) con Semáforo}
\begin{table}
\centering
\small
\begin{tabular}{@{}llcccl@{}}
\toprule
\textbf{KPI} & \textbf{Indicador} & \textbf{Fórmula} & \textbf{Resultado} & \textbf{Meta} & \textbf{Estado} \\
\midrule
KPI$_1$ & $\bar{\Psi}$ Promedio & $\frac{1}{N}\sum \Psi_i$ & \textbf{195.6} & $\leq 250$ & \cellcolor{SmartGreen!25}\semverde{\ding{52} VERDE} \\[6pt]

KPI$_2$ & Error Gypsum & MAPE & \textbf{13.3\%} & $\leq 20\%$ & \cellcolor{SmartGreen!25}\semverde{\ding{52} VERDE} \\[6pt]

KPI$_3$ & Zona Óptima & $\% (100<\Psi\leq 300)$ & \textbf{74.5\%} & $\geq 70\%$ & \cellcolor{SmartGreen!25}\semverde{\ding{52} VERDE} \\[6pt]

KPI$_4$ & $P(\text{riego})_{MC}$ & Montecarlo ($>300$) & \textbf{10.4\%} & $< 50\%$ & \cellcolor{SmartGreen!25}\semverde{\ding{52} VERDE} \\[6pt]

KPI$_5$ & Correlación & $r(\Psi_{ref},\,\Psi_{gyp})$ & \textbf{0.949} & $\geq 0.80$ & \cellcolor{SmartGreen!25}\semverde{\ding{52} VERDE} \\
\bottomrule
\end{tabular}
\end{table}
\end{frame}

\begin{frame}[fragile]{3b. Desglose y Función de Evaluación}
\begin{columns}[c]
\column{0.48\textwidth}
\begin{tcolorbox}[colback=SmartGreen!12,colframe=SmartGreen,
    arc=4pt,left=6pt,right=6pt,top=4pt,bottom=4pt,
    title={\color{white}Distribución Hídrica Real}]
\begin{itemize}
    \item Húmedo ($\leq$100): 10.2\%
    \item \textbf{Óptimo (100--300): 74.5\%}
    \item Seco (300--400): 15.3\%
    \item Crítico ($>$400): \textbf{0.0\%}
\end{itemize}
\end{tcolorbox}

\begin{tcolorbox}[colback=SmartLightBlue!12,colframe=SmartLightBlue,
    arc=4pt,left=6pt,right=6pt,top=4pt,bottom=4pt,
    title={\color{white}Validación del Sensor Gypsum}]
Regresión: $\Psi_{tens} = 0.987\,\Psi_{gyp} + 8.3$\\
$r^2 = 0.901$ \quad $\rightarrow$ \textbf{Alternativa viable}
\end{tcolorbox}

\column{0.48\textwidth}
\begin{tcolorbox}[colback=SmartOrange!12,colframe=SmartOrange,
    arc=4pt,left=6pt,right=6pt,top=4pt,bottom=4pt,
    title={\color{white}Función de Evaluación Python}]
\begin{lstlisting}[numbers=none,frame=none,backgroundcolor={\color{SmartOrange!8}},basicstyle=\ttfamily\small]
def semaforo(val, meta,
             modo='menor'):
    if modo == 'menor':
        ok = val <= meta
    else:
        ok = val >= meta
        
    return 'VERDE' if ok else 'ROJO'
\end{lstlisting}
\end{tcolorbox}
\end{columns}
\end{frame}

% ==================================================================
\section{Programación Orientada a Objetos}
% ==================================================================

% ------------------------------------------------------------------
% DIAPOSITIVA 8: POO - DIAGRAMA DE CLASES
% ------------------------------------------------------------------
\begin{frame}{5a. Arquitectura Orientada a Objetos (POO)}
\begin{columns}[c]
\column{0.5\textwidth}
\begin{tikzpicture}[
    classstyle/.style={rectangle, rounded corners=4pt,
        draw=SmartBlue, fill=SmartBlue!8,
        minimum width=3.5cm, minimum height=0.8cm,
        font=\normalsize\bfseries, text=SmartBlue},
    childstyle/.style={rectangle, rounded corners=4pt,
        draw=SmartGreen!70!black, fill=SmartGreen!10,
        minimum width=3.0cm, minimum height=0.7cm,
        font=\small, text=SmartGray},
    arrow/.style={-Latex, thick, SmartBlue}
]
% Clase base
\node[classstyle] (sensor) at (0,0)
    {\texttt{Sensor (ABC)}};

% Clases hijas
\node[childstyle] (tens) at (-2.0,-2.0)
    {\texttt{Tensiometro}};
\node[childstyle] (resi) at (2.0,-2.0)
    {\texttt{Resistivo}};

% Flechas herencia
\draw[arrow] (tens.north) -- (sensor.south west);
\draw[arrow] (resi.north) -- (sensor.south east);

% Clase Parcela
\node[classstyle, fill=SmartOrange!10, draw=SmartOrange!70] (parcela)
    at (0,-3.8) {\texttt{ParcelaExperimental}};

% Lineas de composicion
\draw[-Latex, dashed, SmartOrange!70] (sensor.south) -- (parcela.north)
    node[midway,font=\footnotesize,right=2pt]{tiene $\rightarrow$};
\end{tikzpicture}

\column{0.45\textwidth}
\begin{tcolorbox}[colback=SmartGreen!8,colframe=SmartGreen!50,
    arc=4pt,left=6pt,right=6pt,top=4pt,bottom=4pt]
\textbf{Abstracción y Herencia:}
\begin{itemize}
    \item Clase base abstracta \texttt{Sensor}.
    \item Las clases hijas heredan los métodos de validación, historial y función \texttt{semaforo()}.
\end{itemize}
\end{tcolorbox}
\end{columns}
\end{frame}

\begin{frame}{5b. Polimorfismo e Instanciación}
\begin{columns}[c]
\column{0.5\textwidth}
\begin{block}{Polimorfismo: método \texttt{leer()}}
\centering
\begin{tabular}{@{}lcc@{}}
\toprule
\textbf{Clase} & \textbf{Entrada} & \textbf{Salida} \\
\midrule
\texttt{Tensiometro} & 300 hPa & 300.0 hPa \\
\multicolumn{3}{l}{\small $\rightarrow$ Lectura directa, sin conversión} \\[4pt]
\texttt{Resistivo} & 300 hPa & 303.5 hPa \\
\multicolumn{3}{l}{\small $\rightarrow$ +Compensación de $T_{suelo}$} \\[4pt]
\texttt{Capacitivo} & $\theta=0.30$ & 98.7 hPa \\
\multicolumn{3}{l}{\small $\rightarrow$ Eq. Van Genuchten inverso} \\
\bottomrule
\end{tabular}
\end{block}

\column{0.45\textwidth}
\begin{tcolorbox}[colback=SmartLightBlue!8,colframe=SmartLightBlue!50,
    arc=4pt,left=6pt,right=6pt,top=4pt,bottom=4pt]
\textbf{Instanciación en la parcela:}\\
\vspace{0.2cm}
Se crearon instancias para los sensores reales del dataset:
\begin{itemize}
    \item \texttt{t42, t43, t44} (Tensiómetro)
    \item \texttt{gyp1, gyp2} (Gypsum)
    \item \texttt{cap1, cap2} (Capacitivo)
    \item \texttt{parcela\_jki} (Parcela)
\end{itemize}
\end{tcolorbox}
\end{columns}
\end{frame}

% ------------------------------------------------------------------
% DIAPOSITIVA 9: POO - CÓDIGO CLAVE
% ------------------------------------------------------------------
\begin{frame}[fragile]{5c. Código POO (Clase Base)}
\begin{exampleblock}{Clase base abstracta}
\begin{lstlisting}[numbers=none,frame=none,basicstyle=\ttfamily\small]
from abc import ABC, abstractmethod

class Sensor(ABC):
    """Clase base para sensores."""
    def __init__(self, id_sensor, ubicacion, profundidad_cm):
        self.id_sensor = id_sensor
        self.ubicacion = ubicacion
        self.historial = []    # Herencia comun
        self.estado    = 'ACTIVO'

    @abstractmethod
    def leer(self, valor_raw, temperatura=None):
        pass  # Implementacion requerida (Polimorfismo)

    def semaforo(self):
        """Mismo algoritmo de evaluacion para TODAS las clases."""
        if not self.historial:
            return 'SIN_DATOS'
        psi = self.historial[-1]['psi']
        if psi <= 100: return 'HUMEDO'
        if psi <= 300: return 'OPTIMO'
        if psi <= 400: return 'SECO'
        return 'CRITICO'
\end{lstlisting}
\end{exampleblock}
\end{frame}

\begin{frame}[fragile]{5d. Código POO (Clase Hija)}
\begin{block}{Clase hija --- SensorResistivo (Gypsum)}
\begin{lstlisting}[numbers=none,frame=none,basicstyle=\ttfamily\small]
class SensorResistivo(Sensor):
    """Gypsum: implementa compensacion termica."""
    TIPO = 'GYPSUM_RESISTIVO'
    COEF_TEMP = 0.035

    def __init__(self, id_sensor, ubicacion, prof_cm,
                 coef_cal=1.0, offset_cal=0.0):
        # Llamada al constructor padre
        super().__init__(id_sensor, ubicacion, prof_cm)
        self.coef_cal   = coef_cal
        self.offset_cal = offset_cal

    def leer(self, valor_raw, temperatura=None):
        # POLIMORFISMO: Logica especifica para Gypsum
        factor = 1.0
        if temperatura is not None:
            factor += self.COEF_TEMP * (temperatura - 20.0)
            
        psi_calc = (valor_raw * self.coef_cal + self.offset_cal) * factor
        return float(max(0, psi_calc))
\end{lstlisting}
\end{block}
\end{frame}

% ==================================================================
\section{Simulación de Montecarlo}
% ==================================================================

% ------------------------------------------------------------------
% DIAPOSITIVA 10: MONTECARLO - CONCEPTO Y CÓDIGO
% ------------------------------------------------------------------
\begin{frame}[fragile]{4a. Simulación de Montecarlo (Código)}
\begin{columns}[c]
\column{0.55\textwidth}
\begin{block}{Función de Simulación}
\begin{lstlisting}[numbers=none,frame=none,basicstyle=\ttfamily\small]
def simulacion_mc_riego(
        psi_media, psi_std,
        n_iter=10000,
        umbral=300.0, semilla=42):
        
    rng = np.random.default_rng(semilla)
    psi_sim = rng.normal(
        loc=psi_media,
        scale=psi_std,
        size=n_iter)
        
    p_riego = np.mean(psi_sim > umbral) * 100
    
    return p_riego
\end{lstlisting}
\end{block}

\column{0.45\textwidth}
\begin{alertblock}{Modelo Estadístico}
$$\Psi_{sim} = \bar{\Psi}_{obs} + \epsilon$$
$$\epsilon \sim \mathcal{N}(0,\sigma_{MC}^2)$$
\vspace{0.3cm}\\
\textbf{Objetivo:}\\
Cuantificar probabilidad de que el estrés hídrico supere el umbral crítico ($300\text{ hPa}$) debido a la incertidumbre.
\end{alertblock}
\end{columns}
\end{frame}

\begin{frame}{4b. Simulación de Montecarlo (Resultados)}
\begin{exampleblock}{Resultados por Escenario (10,000 iteraciones)}
\centering
\small
\begin{tabular}{@{}lccc@{}}
\toprule
\textbf{Escenario temporal} & \textbf{$P$(riego)} & \textbf{IC 95\%} & \textbf{Acción Sugerida} \\
\midrule
Mayo 2016 & \textbf{36.2\%} & [158, 399] & \semamarillo{MONITOREAR} \\[4pt]
Junio 2016 & \textbf{0.2\%} & [26, 246] & \semverde{NO REGAR} \\[4pt]
Julio 2016 & \textbf{0.0\%} & [84, 198] & \semverde{NO REGAR} \\[4pt]
\midrule
\textbf{Período Completo} & \textbf{10.0\%} & \textbf{[16, 362]} & \semverde{\textbf{NO REGAR}} \\
\bottomrule
\end{tabular}
\end{exampleblock}

\vfill

\begin{tcolorbox}[colback=SmartGreen!10,colframe=SmartGreen!50,
    arc=4pt,left=6pt,right=6pt,top=4pt,bottom=4pt]
\centering
\textbf{Convergencia:} El modelo se estabiliza a partir de las $1,000$ iteraciones.
\end{tcolorbox}
\end{frame}

% ==================================================================
\section{Resultados y Visualizaciones}
% ==================================================================

% ------------------------------------------------------------------
% DIAPOSITIVA 11: TABLERO DE CONTROL
% ------------------------------------------------------------------
\begin{frame}{8. Tablero de Control SmartRoot}

\begin{center}
\includegraphics[width=0.88\textwidth,
    height=0.72\textheight,
    keepaspectratio]%
    {../Outputs/figuras/05_tablero_control_completo.png}
\end{center}

\vspace{-4pt}
{\footnotesize\centering
\textbf{Fig. 1:} Panel de control integrado. \textit{Fila superior:} KPI$_1$--KPI$_3$ con semáforo verde.
\textit{Centro:} Serie temporal $\Psi$ con zona óptima sombreada y distribución de estados hídricos.
\textit{Fila inferior:} KPI mensual, Montecarlo del período completo y validación Gypsum--Tensiómetro.
}
\end{frame}

% ------------------------------------------------------------------
% DIAPOSITIVA 12: HISTOGRAMAS MONTECARLO
% ------------------------------------------------------------------
\begin{frame}{8b. Simulación de Montecarlo --- Distribuciones por Escenario}

\begin{center}
\includegraphics[width=0.88\textwidth,
    height=0.68\textheight,
    keepaspectratio]%
    {../Outputs/figuras/04_mc_histogramas_escenarios.png}
\end{center}

\vspace{-4pt}
\begin{columns}[t]
\column{0.5\textwidth}
{\footnotesize
\textbf{Mayo:} distribución centrada en 279~hPa.
Cola derecha $>300$~hPa: $P=36.2\%$ → \semamarillo{MONITOREAR}
}
\column{0.5\textwidth}
{\footnotesize
\textbf{Junio y Julio:} completamente alejados del umbral.
$P(riego) \approx 0\%$ → \semverde{NO REGAR}
}
\end{columns}
\end{frame}

% ------------------------------------------------------------------
% DIAPOSITIVA 13: SERIE TEMPORAL Y CORRELACIONES
% ------------------------------------------------------------------
\begin{frame}{8c. Serie Temporal y Análisis de Correlaciones}

\begin{columns}[t]
\column{0.5\textwidth}
\begin{center}
\includegraphics[width=\textwidth,
    height=0.55\textheight,
    keepaspectratio]%
    {../Outputs/figuras/04_mc_serie_temporal.png}
\end{center}
{\footnotesize \textbf{Convergencia MC:} $\bar{\Psi}_{MC}$ se estabiliza en $N\approx1\,000$ iteraciones.
Panel inferior: $P(\text{riego})$ con clasificación semáforo en tiempo real.}

\column{0.5\textwidth}
\begin{center}
\includegraphics[width=\textwidth,
    height=0.55\textheight,
    keepaspectratio]%
    {../Outputs/figuras/01_correlaciones.png}
\end{center}
{\footnotesize \textbf{Correlación $r=0.949$:} validación del sensor Gypsum.
\textit{Derecha:} Curva de retención Van Genuchten ($\theta_s=0.38$, $\alpha=0.034$, $n=1.37$).}
\end{columns}
\end{frame}

% ------------------------------------------------------------------
% DIAPOSITIVA 14: DIAGRAMA DE FLUJO
% ------------------------------------------------------------------
\begin{frame}{8d. Diagrama de Flujo del Sistema SmartRoot}

\begin{columns}[c]
\column{0.55\textwidth}
\begin{center}
\includegraphics[width=\textwidth,
    height=0.7\textheight,
    keepaspectratio]%
    {../Outputs/figuras/05_diagrama_flujo.png}
\end{center}

\column{0.45\textwidth}
\begin{block}{Pipeline del sistema}
\begin{enumerate}\footnotesize
    \item \textbf{HARDWARE:} Sensor Gypsum (ADC + LoRa)
    \item \textbf{NB01 EDA:} Carga · Validación · NaN
    \item \textbf{NB02 KPI:} Indicadores + semáforo
    \item \textbf{NB03 POO:} Clases Sensor/Parcela
    \item \textbf{NB04 MC:} Montecarlo 10\,000 iter.
    \item \textbf{DECISIÓN:}
        \begin{itemize}\scriptsize
            \item[\semverde{\ding{52}}] $P<20\%$ → No regar
            \item[\semamarillo{$\mathbf{\sim}$}] $20\leq P<50\%$ → Monitorear
            \item[\semrojo{\textbf{!}}] $P\geq 50\%$ → Regar hoy
            \item[\semrojo{\ding{56}}] $P\geq 75\%$ → Urgente
        \end{itemize}
\end{enumerate}
\end{block}

\vspace{4pt}
\begin{tcolorbox}[colback=SmartGreen!8,colframe=SmartGreen!50,
    arc=3pt,left=4pt,right=4pt,top=2pt,bottom=2pt]
\footnotesize \textbf{Resultado del período:}\\
$P(\text{riego})=10\%$ → \semverde{NO REGAR}\\
Cero eventos críticos detectados
\end{tcolorbox}
\end{columns}
\end{frame}

% ==================================================================
\section{Conclusiones}
% ==================================================================

% ------------------------------------------------------------------
% DIAPOSITIVA 15: CONCLUSIONES
% ------------------------------------------------------------------
\begin{frame}{9. Conclusiones}

\begin{columns}[t]
\column{0.5\textwidth}

\begin{exampleblock}{\tick\ Resultados técnicos}
\begin{itemize}\footnotesize
    \item \textbf{5 KPI en verde:} $\bar{\Psi}=195.6$~hPa, MAPE$=13.3\%$, zona óptima$=74.5\%$, $r=0.949$, $P(\text{riego})=10\%$
    \item \textbf{Gypsum validado:} sensor de bajo costo con $r^2=0.90$ respecto al tensiómetro de referencia. \textit{Viable comercialmente}
    \item \textbf{Montecarlo:} mayo 2016 requirió alerta amarilla ($P=36.2\%$); imposible de detectar con umbrales fijos
\end{itemize}
\end{exampleblock}

\begin{block}{\tick\ Valor del Montecarlo}
\footnotesize
La simulación estocástica permite:\\
\textbf{cuantificar el riesgo} en lugar de solo\\
detectarlo, diferencia clave entre\\
sistemas expertos y simples alertas
\end{block}

\column{0.5\textwidth}

\begin{block}{\tick\ Python vs. herramientas estáticas}
\begin{itemize}\footnotesize
    \item \textbf{Reproducibilidad:} semilla fija garantiza resultados idénticos
    \item \textbf{Velocidad:} 2\,494 registros procesados vectorialmente en $<1$~s
    \item \textbf{Escalabilidad:} nuevos sensores = nueva clase hija, sin reescribir lógica central
    \item \textbf{Trazabilidad:} control de versiones Git + notebooks documentados
\end{itemize}
\end{block}

\begin{alertblock}{\tick\ Comprensiones Python}
\footnotesize
Las 3 comprensiones redujeron el código\\
en $\approx60\%$ respecto a ciclos equivalentes,\\
manteniendo legibilidad y eficiencia $O(n)$
\end{alertblock}

\end{columns}
\end{frame}

% ==================================================================
\section{Evolución del Proyecto}
% ==================================================================

% ------------------------------------------------------------------
% DIAPOSITIVA 16: TRABAJO FUTURO
% ------------------------------------------------------------------
\begin{frame}{10. Posible Evolución del Proyecto}

\begin{columns}[t]
\column{0.33\textwidth}
\begin{block}{\small\faRobot\ Machine Learning}
\footnotesize
Reemplazar el modelo estático de Van Genuchten por un \textbf{XGBoost o LSTM} que prediga $\Psi(t+24h)$ usando historial + temperatura + precipitación + ET$_0$\\[4pt]
\textbf{Impacto:} Planificación de riego con 24--72~h de anticipación
\end{block}

\column{0.33\textwidth}
\begin{exampleblock}{\small\faCloud\ IoT + Dashboard Web}
\footnotesize
Conectar SmartRoot a sensores inalámbricos \textbf{ESP32 / LoRaWAN} mediante API MQTT y desplegar el tablero en \textbf{Streamlit} con notificaciones automáticas SMS/email al superar umbral amarillo\\[4pt]
\textbf{Impacto:} Decisiones en tiempo real desde el móvil
\end{exampleblock}

\column{0.33\textwidth}
\begin{alertblock}{\small\faSeedling\ Multi-parcela}
\footnotesize
Implementar clase \texttt{SistemaMultiParcela} que gestione redes distribuidas geoespacialmente, priorizando órdenes de riego por criticidad y distancia logística con \textbf{algoritmos de optimización}\\[4pt]
\textbf{Impacto:} Escala regional / empresarial
\end{alertblock}
\end{columns}

\vspace{8pt}
\begin{tcolorbox}[colback=SmartBlue!8,colframe=SmartBlue!40,
    arc=4pt,left=8pt,right=8pt,top=4pt,bottom=4pt]
\centering\small
\textbf{Visión comercial SmartRoot 2.0:}
Plataforma SaaS de agricultura de precisión con toma de decisiones autónoma,
integrando IA, IoT y simulación estocástica en un único ecosistema Python reproducible
\end{tcolorbox}
\end{frame}

% ------------------------------------------------------------------
% DIAPOSITIVA 17: RESUMEN EJECUTIVO
% ------------------------------------------------------------------
\begin{frame}{Resumen Ejecutivo --- Panel de Resultados}

\begin{columns}[t]
\column{0.5\textwidth}
\begin{block}{Dataset y procesamiento}
\begin{tabular}{@{}ll@{}}
\footnotesize Registros totales & \footnotesize 2\,494 \\
\footnotesize Período & \footnotesize May.--Jul. 2016 \\
\footnotesize Resolución & \footnotesize 30 minutos \\
\footnotesize Sensores modelados & \footnotesize 7 objetos POO \\
\footnotesize Tasa de calidad & \footnotesize $>98\%$ \\
\footnotesize Notebooks & \footnotesize NB01--NB05 \\
\end{tabular}
\end{block}

\vspace{4pt}
\begin{block}{Componentes implementados}
\begin{tabular}{@{}lc@{}}
\footnotesize Programación básica & \footnotesize \tick \\
\footnotesize 3 comprensiones & \footnotesize \tick \\
\footnotesize 3 funciones & \footnotesize \tick \\
\footnotesize 5 KPI con semáforo & \footnotesize \tick \\
\footnotesize POO + herencia + polim. & \footnotesize \tick \\
\footnotesize Montecarlo 10\,000 iter. & \footnotesize \tick \\
\footnotesize Tablero de control & \footnotesize \tick \\
\end{tabular}
\end{block}

\column{0.5\textwidth}
\begin{exampleblock}{KPI globales del período}
\centering
\begin{tabular}{@{}lcc@{}}
\toprule
\footnotesize\textbf{Indicador} & \footnotesize\textbf{Valor} & \footnotesize\textbf{Estado} \\
\midrule
\footnotesize$\bar{\Psi}$ & \footnotesize195.6 hPa & \semverde{\footnotesize VERDE} \\
\footnotesize MAPE & \footnotesize13.3\% & \semverde{\footnotesize VERDE} \\
\footnotesize Zona óptima & \footnotesize74.5\% & \semverde{\footnotesize VERDE} \\
\footnotesize$P$(riego) MC & \footnotesize10.4\% & \semverde{\footnotesize VERDE} \\
\footnotesize$r$ Pearson & \footnotesize0.949 & \semverde{\footnotesize VERDE} \\
\bottomrule
\end{tabular}
\end{exampleblock}

\vspace{4pt}
\begin{tcolorbox}[colback=SmartGreen!12,colframe=SmartGreen,
    arc=4pt,left=6pt,right=6pt,top=4pt,bottom=4pt]
\centering\footnotesize
\textbf{Interpretación ejecutiva}\\
\semverde{Los 5 KPI en estado VERDE}\\
Manejo hídrico \textbf{excelente} en el período\\
$\rightarrow$ SmartRoot: \textbf{cero riegos urgentes} requeridos
\end{tcolorbox}
\end{columns}
\end{frame}

% ------------------------------------------------------------------
% DIAPOSITIVA 18: CIERRE
% ------------------------------------------------------------------
\begin{frame}[plain]
\begin{tikzpicture}[remember picture, overlay]
    \fill[SmartBlue] (current page.north west) rectangle
        ([yshift=-0.45\paperheight]current page.north east);
    \fill[SmartGreen] (current page.south west) rectangle
        ([yshift=0.08\paperheight]current page.south east);
\end{tikzpicture}

\vspace{0.4cm}
\begin{center}
{\color{white}\Huge\bfseries ¡Muchas Gracias!}\\[10pt]
{\color{white!85}\large Proyecto SmartRoot}\\[4pt]
{\color{white!70}\normalsize Sistema Inteligente de Decisión de Riego}\\[4pt]
{\color{white!70}\normalsize Programación Científica en Python}

\vspace{0.6cm}
\begin{tcolorbox}[colback=white!15!SmartBlue,colframe=white!30,
    arc=6pt,left=10pt,right=10pt,top=6pt,bottom=6pt,
    width=0.7\textwidth]
\centering
{\color{white}\large\bfseries ¿Preguntas o comentarios?}\\[6pt]
{\color{white!85}\small
Dataset: Jackisch et al. (2020) · JKI Braunschweig\\
2\,494 registros · Mayo--Julio 2016\\
\textbf{Todos los KPI: \textcolor{SmartGreen!70!white}{VERDE}}}
\end{tcolorbox}

\vspace{0.5cm}
{\color{white!80}\small [Nombre del Estudiante] · Miguel A. Becerra Ph.D.}\\
{\color{white!70}\small Maestría ACI · Fundamentos de Programación · 2026}
\end{center}
\end{frame}

% ============================================================
% DIAPOSITIVAS ADICIONALES (BACKUP / PREGUNTAS)
% ============================================================

\appendix

\begin{frame}[noframenumbering]{Backup: Modelo Van Genuchten implementado}
\begin{columns}[t]
\column{0.5\textwidth}
\begin{block}{Ecuación implementada (clase \texttt{SensorCapacitivo})}

$$\theta(\Psi) = \theta_r + \frac{\theta_s - \theta_r}{\left[1 + (\alpha|\Psi|)^n\right]^{1-1/n}}$$

\footnotesize\textbf{Parámetros suelo JKI:}\\
$\theta_r = 0.050$ · $\theta_s = 0.380$\\
$\alpha = 0.034$~cm$^{-1}$ · $n = 1.370$\\[4pt]
\textbf{Inversión analítica} para convertir\\
$\theta_{medido} \rightarrow \Psi_{estimado}$ (hPa)
\end{block}

\column{0.5\textwidth}
\begin{center}
\includegraphics[width=\textwidth,
    height=0.5\textheight,
    keepaspectratio]%
    {../Outputs/figuras/01_correlaciones.png}
\end{center}
{\footnotesize\centering Panel derecho: curva de retención aplicada}
\end{columns}
\end{frame}

\begin{frame}[noframenumbering]{Backup: Distribución de frecuencias de variables}
\begin{center}
\includegraphics[width=0.88\textwidth,
    height=0.72\textheight,
    keepaspectratio]%
    {../Outputs/figuras/01_histogramas.png}
\end{center}
{\footnotesize\centering Histogramas de $\Psi_{ref}$, $\Psi_{gyp}$, $\theta$ y $T_{suelo}$ con media, mediana y densidad estimada (KDE).}
\end{frame}

\begin{frame}[noframenumbering]{Backup: Patrón de disponibilidad de datos}
\begin{center}
\includegraphics[width=0.88\textwidth,
    height=0.62\textheight,
    keepaspectratio]%
    {../Outputs/figuras/01_patron_nan.png}
\end{center}
\vspace{4pt}
{\footnotesize\centering
\textbf{Verde:} dato disponible. \textbf{Rojo:} dato faltante.
Nota: Gypsum tiene brecha central (registros 1000--3600) detectada automáticamente por SmartRoot.
Sensores capacitivos y de temperatura tienen cobertura $\approx100\%$.}
\end{frame}

\end{document}
